\documentclass[11pt]{article}
\usepackage[margin=1in]{geometry}
\usepackage[T1]{fontenc}
\usepackage{lmodern,amsmath,amssymb,amsthm,mathtools,booktabs,microtype}
\usepackage[colorlinks=true,linkcolor=blue!55!black,urlcolor=blue!55!black,citecolor=blue!55!black]{hyperref}
\usepackage{xcolor}
\newtheorem{theorem}{Theorem}[section]
\newtheorem{lemma}[theorem]{Lemma}
\newtheorem{corollary}[theorem]{Corollary}
\newtheorem{proposition}[theorem]{Proposition}
\theoremstyle{remark}\newtheorem{remark}[theorem]{Remark}
\newcommand{\F}{\mathbb F}
\newcommand{\E}{\mathbb E}
\newcommand{\Z}{\mathbb Z}
\newcommand{\C}{\mathbb C}
\newcommand{\1}{\mathbf 1}
\newcommand{\norm}[1]{\left\lVert#1\right\rVert}
\newcommand{\abs}[1]{\left|#1\right|}
\newcommand{\spt}{\operatorname{supp}}
\newcommand{\codim}{\operatorname{codim}}
\newcommand{\ord}{\operatorname{ord}}
\title{\textbf{Sharp Euler-Factor Separation for Products of Odd-Prime-Order Dirichlet $L$-Functions}\\[4pt]\large Projection energy, extremal cosets, and robust prime certificates}
\author{Research note for the ProveIt GRH project}
\date{27 November 2026}
\begin{document}
\maketitle
\begin{abstract}
Let $X$ be an elementary abelian $\ell$-group of Dirichlet characters, $\ell$ prime, of rank $r\ge1$. For distinct multisets of $k$ characters, the density $\delta$ of primes at which their local Euler polynomials differ satisfies the sharp universal inequality $\delta\ge(\ell-1)/(\ell k)$. Equality occurs precisely for two distinct parallel cosets of a subspace of size $k=\ell^d$; thus it is completely classified in every rank. For $k\le\ell$ and $r\ge2$ we obtain the stronger independent bound $\delta\ge(\ell-1)(\ell+1-k)/\ell^2$, classify its equality geometry, and attain it for $k\mid(\ell-1)$, $k\mid(\ell+1)$ with $k<\ell+1$, and $k=\ell$. We also determine the exact packet-size thresholds for the two smallest positive disagreement densities in every rank. Under Dirichlet GRH, an energy minorant of Fourier algebra norm at most two transfers separation to primes in $(x/2,x)$ at the uniform scale $x\gg k^2\log^2(2Q)$, with robust decoding against adversarially corrupted Euler factors. The density theorems themselves are unconditional. These results partially answer an explicit open question in the ProveIt merged article.
\end{abstract}
\section{Source and scope}
This note extends the article's subsection \emph{Matching Euler factors of products of quadratic $L$-functions} and its open question about odd prime order \cite{article}. The article's theorem \texttt{rec:thm:qmatch} uses Walsh characters and dyadic valuations to classify minimum and near-minimum disagreement for quadratic packets. Its theorem \texttt{rec:thm:qmatchfinite} transfers Fourier-algebra estimates to short multiplicative prime windows. In the odd-order setting, the right observable is not a Walsh sign but a whole histogram of a linear projection; we exploit the integral projection energy, which controls the Fourier algebra norm uniformly in the exponent, rank and packet size.

The vendored Lean development \cite{vendor} identifies the upstream revision as\\ \texttt{72fead102f663d5c7a407964a13dc0cf9d977207}. The inspected statement file defines\\ \texttt{DirichletLNonvanishingRightOfHalf} and \texttt{DirichletLZerosOnLine}. The former quantifies over all moduli $N\ge1$, all Dirichlet characters modulo $N$, and all complex $s$ with $\Re s>1/2$, excluding the principal-character pole point $s=1$; the latter locates every zero in $0<\Re s<1$ on $\Re s=1/2$. The inspected \texttt{src/Main.lean} exports the corresponding theorems \texttt{LFunction\_ne\_zero\_of\_half\_lt\_re} and \texttt{LFunction\_eq\_zero\_re}.\par The first is obtained from \texttt{New.grh}, proved in the file \texttt{OneSidedBound.lean} in the reduction directory. The upstream README reports independent proof checks and the three usual Lean axioms. We have inspected these interfaces, not rebuilt the vendored project or independently checked its complete proof. The new results are classical written proofs, not Lean-certified. The finite-prime conclusions use the stated Dirichlet GRH theorem as input; no GRH for general Hecke, Artin, or automorphic functions is needed.

\section{The exact finite-geometric model}
Fix a modulus $Q\ge 2$, a prime $\ell$, and a subgroup $X\subseteq\widehat{(\Z/Q\Z)^\times}$ of exponent $\ell$ and rank $r\ge1$. Write $X$ additively as $\F_\ell^r$ (addition corresponds to character multiplication) and $V=\operatorname{Hom}_{\F_\ell}(X,\F_\ell)$. Fix a primitive $\ell$th root of unity $\zeta$. For each unit $u$ modulo $Q$, evaluation defines $v_u\in V$ by $\chi_a(u)=\zeta^{v_u(a)}$. The map $u\mapsto v_u$ is onto: otherwise the annihilator of its image would contain a nonzero character which evaluates to one on every unit. Its fibers have equal cardinality.

A \emph{packet} $A$ is a multiset of $k$ elements of $X$. At a rational prime $p\nmid Q$, its inverse Euler factor is the polynomial
\[
 P_A(p,T)=\prod_{a\in A}(1-\chi_a(p)T).
\]
Let $A,B$ be distinct packets of equal size $k\ge1$, and set
\[
 c(a)=\operatorname{mult}_A(a)-\operatorname{mult}_B(a),\qquad
 s=\sum_{c(a)>0}c(a)=\tfrac12\sum_a|c(a)|\in[1,k].
\]
We refer to the two disjoint residual multisets of size $s$ as the cancelled packets. For $v\in V$ and $t\in\F_\ell$ put
\[
 d_v(t)=\sum_{v(a)=t}c(a),\quad E(v)=\sum_{t\in\F_\ell}d_v(t)^2,
 \quad H(v)=\sum_{a\in X}c(a)\zeta^{v(a)}.
\]
\begin{lemma}[Local factors, projections and traces]\label{lem:projection}
The following conditions are equivalent: the local polynomials at a unit with evaluation $v$ differ; $d_v$ is not identically zero; $E(v)>0$; $H(v)\ne0$. The disagreement set $D\subset V$ is a union of punctured one-dimensional subspaces, excludes $0$, and is nonempty. Consequently the prime disagreement density exists unconditionally and is
\[
 \delta(A,B)=\frac{|D|}{\ell^r}=\frac{(\ell-1)J}{\ell^r},
 \qquad 1\le J\le\frac{\ell^r-1}{\ell-1}.
\]
Here $J$ is the number of projective disagreement directions.
\end{lemma}
\begin{proof}
Unique factorization of the polynomials gives equality exactly when all projected multiplicities agree. This is equivalent to $d_v=0$, hence to $E(v)=0$. If $H(v)=0$, the integral polynomial $\sum_{t=0}^{\ell-1}d_v(t)Z^t$ is divisible over $\mathbb Q$ by $\Phi_\ell(Z)=1+Z+\cdots+Z^{\ell-1}$. Its degree is at most $\ell-1$ and its value at $1$ is zero, so it is identically zero. The reverse implication is immediate. Multiplying $v$ by a nonzero scalar permutes the bins, and $d_0=0$. If $D$ were empty then all Fourier coefficients of $c$ would vanish, contrary to $c\ne0$. The prime number theorem for arithmetic progressions and the equal fibers of the evaluation map give the density, with finitely many primes $p\mid Q$ immaterial.
\end{proof}

This trace equivalence uses \emph{prime} exponent and equal degree; it does not hold in this form for arbitrary composite exponents. The entire paper permits the trivial character as a member of $X$; it is naturally included in a group of exponent $\ell$.

\section{Universal bound and all equality cases}
\begin{theorem}[Sharp projection-energy separation]\label{thm:energy}
For every pair of distinct size-$k$ packets,
\begin{equation}\label{eq:energybound}
 \delta\ \ge\ \frac{\ell-1}{\ell}\frac{\sum_a c(a)^2}{2s^2}
 \ \ge\ \frac{\ell-1}{\ell s}\ \ge\ \frac{\ell-1}{\ell k}.
\end{equation}
Equality in the last universal bound holds if and only if $k=\ell^d$ for an integer $0\le d\le r-1$ and, up to swapping packets, $A=a+U$, $B=b+U$ are disjoint full cosets of the same $d$-dimensional subspace $U\subseteq X$. In particular, there are no repeated members or common members at equality. Thus if $k=\ell^d$ with $d<r$, the exact minimum over all packets is $(\ell-1)/\ell^{d+1}$.
\end{theorem}
\begin{proof}
Expand the squares of the projections:
\[
 E(v)=\sum_{a,b}c(a)c(b)\1_{v(a-b)=0}.
\]
For uniform $v$, the expectation of the indicator is $1$ when $a=b$ and $1/\ell$ otherwise. Since $\sum c=0$,
\begin{equation}\label{eq:identity}
 \E_{v\in V} E(v)=(1-1/\ell)\sum_a c(a)^2.
\end{equation}
For each $v$, the positive and negative masses of $d_v$ are equal to some $s_v\le s$. Therefore
\[
 E(v)\le 2s_v^2\le2s^2,
 \qquad E(v)/(2s^2)\le\1_D(v).
\]
The bound follows by averaging and using integrality $\sum c(a)^2\ge\sum|c(a)|=2s$.

For equality with $(\ell-1)/(\ell k)$, all three inequalities must be equalities: $s=k$, every nonzero coefficient is $\pm1$, and $E(v)=2k^2$ for every $v\in D$. The latter happens exactly when each packet projects to a single bin, with the two bins distinct. Indeed equality in the bound for a sum of squares of nonnegative masses forces only one positive bin and one negative bin; $s_v=s$ excludes any projected cancellation.

Choose $a\in A,b\in B$ and let $U$ be the span of all within-packet differences, from both packets. Then $A\subseteq a+U$, $B\subseteq b+U$. For a disagreeing $v$ we have $v\in U^\perp$ and $v(a-b)\ne0$. Conversely these two conditions give distinct constant projections and hence disagreement. Thus
\[
 D=U^\perp\setminus(U^\perp\cap(a-b)^\perp).
\]
It is nonempty, so $a-b\notin U$. Writing $d=\dim U$, we have $\delta=(\ell-1)/\ell^{d+1}$. Equality with the asserted universal bound forces $k=\ell^d$. Both size-$k$ sets must therefore fill their cosets. The same formula for $D$ verifies the converse for every pair of distinct full cosets.
\end{proof}
\begin{corollary}[Quantization and rank-one case]\label{cor:rounding}
For arbitrary $k$, $J\ge\lceil\ell^{r-1}/k\rceil$. For $r=1$, every distinct pair of equal-size packets disagrees at density $(\ell-1)/\ell$.
\end{corollary}
\begin{proof}
Combine Lemma~\ref{lem:projection} with Theorem~\ref{thm:energy}. In rank one there is exactly one projective direction and it must disagree.
\end{proof}
\begin{remark}[A useful exact deficit]
Set $F(v)=E(v)/(2s^2)$. The identity
\[
 \delta-\frac{\ell-1}{\ell s}
 =\E(\1_D-F)+\frac{\ell-1}{2\ell s^2}\sum_a (c(a)^2-|c(a)|)
\]
separates two obstructions to extremality: nonconstant projections on disagreeing directions, and repeated residual characters. Both terms are nonnegative. In particular, if the left side is at most $\varepsilon(\ell-1)/(\ell s)$, then $\sum_a(c(a)^2-|c(a)|)\le2\varepsilon s$ and $\E(\1_D-F)\le\varepsilon(\ell-1)/(\ell s)$. This is a quantitative, though not yet geometric, stability statement.
\end{remark}

\section{Small packets: a stronger bound from hyperplane pencils}
\begin{theorem}[Planar extremizers]\label{thm:pencil}
Suppose $r\ge2$ and $1\le k\le\ell$. Then
\begin{equation}\label{eq:pencil}
 \delta\ge b_{\ell,k}:=\frac{(\ell-1)(\ell+1-k)}{\ell^2}.
\end{equation}
For $2\le k\le\ell$, equality is equivalent to the following structure: $A,B$ are disjoint $k$-element sets in a common affine plane of $X$; the cross differences $a-b$ lie in exactly $k$ one-dimensional directions in that plane, and each direction occurs exactly once in the cross differences incident with each vertex of either set. In this case the agreement set in $V$ is the union of the $k$ distinct hyperplanes with these normals, all through the same codimension-two subspace. For $k=1$, every pair attains the bound.
\end{theorem}
\begin{proof}
Choose an element $a$ of the cancelled positive packet. If the two projected multisets agree at $v$, at least one of the $s$ negative members $b$ satisfies $v(a-b)=0$. Hence the agreement set lies in the union of at most $s\le k$ distinct linear hyperplanes. Write their actual number as $t$. They are genuine hyperplanes since the cancelled packets are disjoint. For $t\ge1$,
\[
 \left|\bigcup_{i=1}^{t}H_i\right|
 \le \ell^{r-1}+(t-1)(\ell^{r-1}-\ell^{r-2}),
\]
by counting $H_1$ and each $H_i\setminus H_1$ for $i>1$. Subtract from $\ell^r$, divide by $\ell^r$, and use $t\le k$. If $t\ge2$, equality in the union bound requires the sets $H_i\setminus H_1$ ($i>1$) to be disjoint. For $t=2$ the hyperplanes already share a codimension-two intersection. For $t\ge3$, $H_i\cap H_j\subseteq H_1$ for $i,j>1$, so, by dimension, $H_i\cap H_j=H_i\cap H_1=H_j\cap H_1$. Conversely a common codimension-two intersection makes the outside pieces disjoint.

At equality in \eqref{eq:pencil}, for every positive anchor $a$ there must be $k$ distinct hyperplanes and $s=k$. Thus the negative packet has $k$ distinct members. Reversing the roles gives the same for the positive packet. All these pencils have exactly the agreement set as their union. Since $k\le\ell$, no other hyperplane can be contained in this union: any other hyperplane meets each of the $k$ in a codimension-one subspace of itself. If $k<\ell$, their sizes sum to less than its size, and if $k=\ell$ their common origin makes the sum a strict overcount (in rank two the intersections are just the origin). Thus the set of $k$ hyperplanes is the same for every anchor. Their normals span a two-dimensional $U\subset X$. Every $a-b$ lies in $U$, so both sets lie in one affine coset of $U$. For each anchor the $k$ directions occur once. Conversely, for any of these directions, the cross differences along it give a perfect matching of the two sets, which projects identically on its annihilating hyperplane. No other $v$ can give equal projections by the anchor argument. This proves the equality characterization.
\end{proof}

\begin{theorem}[Two infinite families of sharp constructions]\label{thm:constructions}
For every prime $\ell$ and rank $r\ge2$, the bound \eqref{eq:pencil} is attained if $k\mid(\ell-1)$, or $k\mid(\ell+1)$ and $k<\ell+1$ with $k\le\ell$, or $k=\ell$. Thus it is the exact minimum for all these sizes, and in particular for $k=1,2,3$ whenever $k\le\ell$.
\end{theorem}
\begin{proof}
All constructions use a two-dimensional subspace, and restriction of $v$ to it is uniform. For $k\mid(\ell-1)$, take the multiplicative subgroup $M\subset\F_\ell^\times$ of size $k$ and the sets
\[
 A=\{(m,0):m\in M\},\qquad B=\{(0,m):m\in M\}.
\]
Their projections under $(u,v)$ agree at the origin and when $u,v\ne0$ with $u/v\in M$. These are exactly $k$ projective directions, so the number of nonzero disagreement directions is $\ell+1-k$. The claimed density follows.

For the nonsplit construction, view the plane as $\F_{\ell^2}$, take a subgroup $M$ of size $k$ of the cyclic norm-one group $\{m:m^{\ell+1}=1\}$, and choose $z\ne0$ with $z^{\ell-1}\notin M$ (the map $z\mapsto z^{\ell-1}$ is onto the norm-one group). Then $A=Mz$ and $B=Mz^\ell$ are disjoint. For $m\in M$, the $\F_\ell$-linear involution $R_m(w)=m w^\ell$ swaps $A$ and $B$. When $\ell$ is odd, it is not scalar and has a one-dimensional fixed space on the dual plane. For distinct $m,n$, the product $R_mR_n$ is multiplication by $m/n\ne1$, with no nonzero fixed dual vectors, so these fixed dual lines are distinct. Each is an agreement direction. Theorem~\ref{thm:pencil} shows there are no additional directions. The only case $\ell=2$ with a proper divisor $k\le2$ of $\ell+1=3$ is $k=1$, already settled. Finally for $k=\ell$, take two distinct parallel affine lines; Theorem~\ref{thm:energy} gives $\delta=(\ell-1)/\ell^2=b_{\ell,\ell}$. For $k=3<\ell$, the prime $\ell$ is congruent to $\pm1$ modulo $3$; all other small cases are immediate.
\end{proof}
\begin{remark}[Polynomial certificate for planar equality]
For planar $A,B$ of size $k\le\ell$, put $R_A(T,u,v)=\prod_{a\in A}(T-u a_1-v a_2)$, and similarly $R_B$. If there are $k$ agreement directions, every homogeneous coefficient difference of degree less than $k$ vanishes identically: a binary homogeneous polynomial of degree $j$ vanishing along $k>j$ distinct lines is zero. The remaining constant-in-$T$ difference of degree $k$ is nonzero and is a scalar times the product of the $k$ defining linear forms of the agreement lines. Conversely such an identity implies agreement at those lines, with no further agreement by the bound. This offers an exact finite-field search certificate for the remaining sizes. It is a statement about polynomials over $\F_\ell$, not an approximation of complex roots of unity.
\end{remark}

\section{The two highest-agreement levels in every rank}
\begin{theorem}[Exact thresholds]\label{thm:levels}
For $r\ge2$, the smallest two positive density levels allowed by projective quantization, $J=1,2$, are attainable at packet size $k$ precisely under the respective conditions
\[
 J=1:\ k\ge\ell^{r-1},\qquad
 J=2:\ k\ge(\ell-1)\ell^{r-2}.
\]
In particular the exact minimum disagreement density for $k\ge(\ell-1)\ell^{r-2}$ is
\[
 \delta_{\min}(k)=
 \begin{cases}
  2(\ell-1)/\ell^r,&(\ell-1)\ell^{r-2}\le k<\ell^{r-1},\\
  (\ell-1)/\ell^r,&k\ge\ell^{r-1}.
 \end{cases}
\]
\end{theorem}
\begin{proof}
If $J=1$, Fourier inversion on $X$ expresses $c(a)$ as an integral-valued function of a single nonzero linear form $v_1(a)$. Its mean is zero and it is nonzero; each positive bin represents $\ell^{r-1}$ elements, so $s\ge\ell^{r-1}$. Two distinct parallel hyperplanes attain this value by Theorem~\ref{thm:energy}.

If $J=2$, choose nonzero representatives $v_1,v_2$ of the two directions, necessarily linearly independent. Fourier inversion gives
\[
 c(a)=f(v_1(a),v_2(a)),\qquad f(i,j)=u(i)+w(j),
\]
where both univariate functions are nonconstant (otherwise one Fourier direction would be absent). The real-valued representation can be made integral by using $u(i)=f(i,0)$, $w(j)=f(0,j)-f(0,0)$, since the map $a\mapsto(v_1(a),v_2(a))$ is onto. Each pair $(i,j)$ has $\ell^{r-2}$ preimages. We claim at least $2(\ell-1)$ entries of the $\ell$ by $\ell$ matrix $f$ are nonzero. If no value of $u$ occurs more than $\ell-2$ times, every column has at least two nonzero entries. Otherwise, since $u$ is nonconstant, some value $u_0$ occurs $\ell-1$ times. Every column has at least one nonzero entry, and nonconstancy of $w$ ensures at least one column with $w(j)\ne-u_0$; that column has at least $\ell-1$ nonzero entries. This proves the claim in either case, including $\ell=2$. Because the nonzero entries are integers and the total sum of $c$ is zero, we get $s=\norm c_1/2\ge(\ell-1)\ell^{r-2}$.

For the converse take $c(a)=\1_{v_1(a)=0}-\1_{v_2(a)=0}$. Its mean is zero; its Fourier support is exactly the two punctured lines generated by $v_1,v_2$. Its positive mass is $(\ell-1)\ell^{r-2}$, as counted from the two hyperplanes after deleting their intersection. In both constructions pad the packets with the same arbitrary characters to reach any larger $k$. The minimum formula follows from quantization and the thresholds.
\end{proof}

\section{GRH transfer: separation by short primes and robust decoding}
Choose once and for all a nonnegative smooth weight $w$ supported in $(1/2,1)$ with $\int w=1$. Define
\[
 M_x(f)=\sum_{p\nmid Q}(\log p)w(p/x)f(v_p),\qquad
 W_x=M_x(1),\quad \widehat f(a)=\E_{v\in V} f(v)\zeta^{-v(a)},\quad
 \norm f_{\mathcal A}=\sum_a|\widehat f(a)|.
\]
\begin{lemma}[Smoothed prime transfer; uses Dirichlet GRH]\label{lem:transfer}
There exist constants $C_w,x_w$ depending only on $w$ such that, for $x\ge x_w$ and every $Q,X,f$ as above,
\[
 \left|\frac{M_x(f)}x-\E f\right|\le e_x\norm f_{\mathcal A},
 \qquad e_x=C_w\frac{\log(2Q)}{\sqrt{x}}.
\]
\end{lemma}
\begin{proof}
Apply the article's prime Gram estimate \texttt{rec:lem:gram} (proved by the smoothed explicit formula under Dirichlet GRH) to a primitive character inducing each member $a\in X$, paired with the primitive principal character of conductor one. Each conductor is at most $Q$; removing extra primes dividing $Q$ costs $O_w(\log Q)$. In the principal case the main term is $x\int w=x$. Expand $f$ in its characters of $V$ and use the triangle inequality.

For clarity, the analytic estimate uses Mellin inversion of $-L'/L$, the pole at one for the principal character, the absolutely summed zero contribution $O_w(\sqrt{x}\log(2Q))$ from $\Re\rho=1/2$, and the bound $O_w(\sqrt{x})$ for prime powers. The rapidly decaying Mellin transform of $w$ and the standard local zero count make the implied constant uniform in the conductor. The full proof of this input is the cited article lemma; see also \cite{davenport,mv}. This is the only GRH input to the present argument.
\end{proof}
\begin{theorem}[Uniform finite-prime separation and decoding; uses Dirichlet GRH]\label{thm:finite}
Let $d_0=(\ell-1)/(\ell k)$. Then, simultaneously for all distinct size-$k$ packets,
\begin{equation}\label{eq:finite}
 \frac{M_x(\1_D)}{x}\ge
 \frac{\ell-1}{\ell}\frac{\sum_a c(a)^2}{2s^2}-2e_x\ge d_0-2e_x.
\end{equation}
Whenever $W_x>0$ and $e_x<1$, their normalized weighted Hamming distance among primes in $(x/2,x)$ is at least $(d_0-2e_x)/(1+e_x)$ provided the numerator is nonnegative. Given $0<\varepsilon<1$, at
\[
 x\ge\max\{x_w, C'_w\varepsilon^{-2} k^2\log^2(2Q)\}
\]
this distance is at least $(1-\varepsilon)d_0$. There is therefore a distinguishing prime in this interval. Any list of Euler factors of an unknown packet of size $k$ that is corrupted at primes of total normalized weight strictly less than $(1-\varepsilon)d_0/2$ uniquely identifies the packet by minimum weighted Hamming distance. More generally, if erasures have weight $u$ and erroneous unerased coordinates weight $v$, then $u+2v<(1-\varepsilon)d_0$ suffices.
\end{theorem}
\begin{proof}
The energy minorant is $F=E/(2s^2)$, so $0\le F\le\1_D$. For $a\ne b$, the identity\[\1_{v(a-b)=0}=\ell^{-1}\sum_{j\in\F_\ell}\zeta^{v(j(a-b))}\] shows that the Fourier algebra norm is one; for $a=b$ the function is constant one. Expanding $E$ as above gives
\[
 \norm F_{\mathcal A}\le \frac{(\sum_a|c(a)|)^2}{2s^2}=2.
\]
Now apply Lemma~\ref{lem:transfer} and the exact identity \eqref{eq:identity}. This proves \eqref{eq:finite}. Also $|W_x/x-1|\le e_x$. For example choose $C'_w=64C_w^2$, so $e_x\le\varepsilon/(8k)\le\varepsilon d_0/4$. Then $W_x>0$ and
\[
 \frac{d_0-2e_x}{1+e_x}\ge d_0-3e_x\ge(1-\varepsilon)d_0.
\]
The weighted Hamming distance is a metric on the distinct packet codewords (all zero-weight coordinates are ignored). Triangle inequality shows that two distinct codewords cannot each agree with a received word up to errors of weight $v$ and erasures of weight $u$ when $u+2v$ is smaller than the distance lower bound. A finite exhaustive search over the $\binom{|X|+k-1}{k}$ packets is one decoder; no efficient decoding complexity is asserted.
\end{proof}
\begin{corollary}[Sharp extremal density measurements; uses Dirichlet GRH]\label{cor:measure}
For any extremal coset pair of Theorem~\ref{thm:energy}, $\norm{\1_D}_{\mathcal A}=2k\delta$. The same formula holds for every pencil equality case of Theorem~\ref{thm:pencil}. In either case,
\[
 |M_x(\1_D)/x-\delta|\le 2k\delta\,e_x.
\]
In particular the density can be measured to relative precision $\varepsilon$ already at $x\gg_w\varepsilon^{-2}k^2\log^2(2Q)$, not a scale growing with the rank. The same relative accuracy (with a different absolute constant) holds with normalization by $W_x$.
\end{corollary}
\begin{proof}
For cosets $D=K\setminus K'$ with $K=U^\perp$, $K'=K\cap(a-b)^\perp$ and $[K:K']=\ell$. The Fourier coefficient of a subgroup indicator is its relative size on its annihilator, zero elsewhere. Hence $\norm{\1_K-\1_{K'}}_{\mathcal A}=2(1-1/\ell)=2k\delta$.

For a pencil of $k\ge2$ hyperplanes $H_i$ with common codimension-two $W$, use
\[
 \1_D=1-\sum_i\1_{H_i}+(k-1)\1_W.
\]
Its coefficient at zero is $\delta$. On each of the $k$ normal lines, at nonzero elements, its coefficient is $-(\ell+1-k)/\ell^2$; on the remaining $\ell+1-k$ lines of the normal plane it is $(k-1)/\ell^2$; all other coefficients vanish. The absolute sum is $\delta+(\ell-1)[k(\ell+1-k)+(\ell+1-k)(k-1)]/\ell^2=2k\delta$. For $k=1$ the coset calculation applies. Lemma~\ref{lem:transfer} gives the estimate. Normalization changes relative error by at most $(2k+1)e_x/(1-e_x)$.
\end{proof}
\begin{remark}[A stronger finite bound for small packets]
In the proof of Theorem~\ref{thm:pencil}, the signed function $1-\1_{H_1}-\sum_{i>1}(\1_{H_i}-\1_{H_i\cap H_1})$ is a pointwise minorant of $\1_D$, with mean at least $b_{\ell,k}$ and Fourier algebra norm at most $2k$. Thus for $k\le\ell$, $M_x(\1_D)/x\ge b_{\ell,k}-2ke_x$ as well as \eqref{eq:finite}. This need not be positive at the same scale near $k=\ell$; the uniform energy theorem avoids that loss. In the exact pencil case Corollary~\ref{cor:measure} removes it altogether.
\end{remark}

\section{Examples, significance and remaining questions}
\begin{center}
\begin{tabular}{rrrrl}
\toprule
$\ell$ & $r$ & $k$ & $\delta_{\min}$ & justification\\
\midrule
3 & $\ge2$ & 1 & $2/3$ & single characters\\
3 & $\ge2$ & 2 & $4/9$ & split pencil\\
3 & $\ge2$ & 3 & $2/9$ & parallel cosets\\
3 & 3 & 6,7,8 & $4/27$ & two-line threshold\\
3 & $\ge3$ & 9 & $2/27$ & parallel cosets\\
5 & $\ge2$ & 3 & $12/25$ & nonsplit pencil\\
5 & $\ge2$ & 4 & $8/25$ & split pencil\\
7 & $\ge2$ & 4 & $24/49$ & nonsplit pencil\\
7 & $\ge3$ & 49 & $6/343$ & parallel cosets\\
\bottomrule
\end{tabular}
\end{center}
In particular two different products of nine cubic Dirichlet factors can agree at $25/27$ of all primes, and no such pair can agree more often. All extremizers are exactly two parallel affine planes of characters. Nevertheless, an effective prime pool of scale $O(k^2\log^2 Q)$ with fixed relative margin distinguishes them and all other size-$k$ products. The construction is available for every abstract rank: choose $r$ distinct rational primes congruent to one modulo $\ell$ and take $Q$ to be their product. Each cyclic unit group has a quotient of order $\ell$, yielding independent order-$\ell$ characters by the Chinese remainder theorem.

The unconditional result is not merely the ordinary uncertainty principle $\delta\ge1/(2k)$: the integral signed projection energy improves the constant to $(1-1/\ell)/k$ for every odd prime exponent and characterizes equality. Meshulam's refined uncertainty inequality \cite{meshulam} already yields the numerical lower bound at $k=\ell^d$ (apply it between consecutive powers of $\ell$ to Fourier support size at most $2k$); we do not claim that special lower bound as new in isolation. The stronger all-size integral bound, equality analysis, planar characterization and constructions, two-level thresholds, and their uniform finite-prime transfer form the extension here. The projection identity belongs to the familiar finite Radon/Fourier toolkit, and equal projections are a standard discrete tomography problem \cite{tomography,fraser}. For $\ell=2$ our all-size inequality need not recover the sharper dyadic staircase in the article; the intended advance is the odd-prime case. Thus the novelty claim is relative to the inspected article and to the specific combination of conclusions; the limited literature search cannot certify publication priority. No famous conjecture outside the article is claimed solved.

\paragraph{Further research questions.}
(1) Determine exactly which $k\le\ell$ admit the equality geometry of Theorem~\ref{thm:pencil}; the split and nonsplit constructions cover many but not all sizes. (2) Determine the minimum positive mass for $J=3,4,\ldots$ projective directions, including their dependence on the dimension of the span of the directions. (3) Prove quantitative geometric stability near the extremizers of Theorem~\ref{thm:energy}, strengthening the exact deficit identity to a distance from parallel cosets. (4) Seek composite-exponent analogues using simultaneous traces at several powers and the full cyclotomic factorization. (5) Formalize the finite algebra lemmas in Lean, and connect a formal smoothed prime estimate to the pinned GRH interfaces. A formalization roadmap, exact finite checks and source audit accompany the note.

\begin{thebibliography}{9}
\bibitem{article} ProveIt, \emph{Beyond the Riemann hypothesis for Dirichlet $L$-functions: consequences, transfers and limits}, merged article, sections \texttt{src.tex} and \texttt{rec.tex}, in particular \texttt{rec:lem:gram}, \texttt{rec:thm:qmatch}, \texttt{rec:thm:qmatchfinite} and \texttt{rec:sec:open}. \href{https://github.com/VladimirReshetnikov/ProveIt/tree/main/Analysis/GeneralizedRiemannHypothesis/Article}{Repository directory}, accessed 27 Nov.
\bibitem{vendor} \texttt{lsernau/grh-lean}, \emph{The Generalized Riemann Hypothesis for Dirichlet L-functions, in Lean 4}, readable edition, upstream revision \texttt{72fead102f663d5c7a407964a13dc0cf9d977207}. \href{https://github.com/VladimirReshetnikov/ProveIt/tree/main/Analysis/GeneralizedRiemannHypothesis/vendor/grh-lean}{Vendored source}. Proof attributed there to Levent Alp\"oge with Claude.
\bibitem{davenport} H. Davenport, \emph{Multiplicative Number Theory}, 3rd ed., Springer, 2000.
\bibitem{mv} H. L. Montgomery and R. C. Vaughan, \emph{Multiplicative Number Theory I: Classical Theory}, Cambridge University Press, 2007.
\bibitem{meshulam} R. Meshulam, An uncertainty inequality for finite abelian groups, \emph{European J. Combin.} \textbf{27} (2006), 63--67. \href{https://doi.org/10.1016/j.ejc.2004.07.009}{DOI}.
\bibitem{tomography} J. Matou\v{s}ek, A. P\v{r}\'\i v\v{e}tiv\'y and P. \v{S}kro\v{n}, How many points can be reconstructed from $k$ projections?, \emph{SIAM J. Discrete Math.} (2008). \href{https://doi.org/10.1137/080715706}{DOI}.
\bibitem{fraser} J. M. Fraser and F. Rakhmonov, Exceptional projections in finite fields: Fourier analytic bounds and incidence geometry, \href{https://arxiv.org/abs/2503.15072}{arXiv:2503.15072} (2025).
\end{thebibliography}
\end{document}
